\documentclass[12pt,letterpaper]{article}

\usepackage[margin=1in]{geometry}
\usepackage{amsmath,amssymb}
\usepackage{fancyhdr}

\setlength{\parindent}{0pt}
\setlength{\parskip}{0.8em}
\setlength{\headheight}{15pt}

\pagestyle{fancy}
\fancyhf{}
\fancyhead[L]{FIT3155 SAMPLE EXAM PAPER}
\fancyfoot[C]{Page \thepage{} of 22}
\renewcommand{\headrulewidth}{0pt}
\renewcommand{\footrulewidth}{0pt}

\newcommand{\marksbox}[1]{%
  \vfill
  \begin{flushright}
    \fbox{\parbox[c][1.1cm][c]{1.3cm}{\centering #1}}
  \end{flushright}
}

\newcommand{\questiontitle}[2]{%
  \textbf{Question #1}\hfill\textbf{(#2 marks)}
}

\newcommand{\workingpage}{%
  \newpage
  \begin{center}
    (Page intentionally left blank for working space)
  \end{center}
}

\begin{document}

\begin{titlepage}
  \thispagestyle{empty}
  \vspace*{0.18\textheight}
  {\Large FIT3155 \hfill SAMPLE EXAM PAPER\par}
  \vspace{2.2cm}
  \begin{center}
    {\LARGE\bfseries FIT3155 SAMPLE EXAM PAPER\par}
    \vspace{0.8cm}
    {\large Generated Sample Exam 03\par}
  \end{center}
\end{titlepage}

\setcounter{page}{2}

\section*{INSTRUCTIONS}

\begin{itemize}
  \item This exam is worth 60 marks.
\end{itemize}

\begin{center}
  \textbf{General exam technique}
\end{center}

\begin{itemize}
  \item Some students lose marks by not attempting all questions.

  \item Suppose you are likely to get 3/6 on a question after spending say 10 minutes on it.
  Spending another 10 minutes on the same question may gain you at most 3 additional
  marks. In contrast, spending those 10 minutes on a new question could earn you another
  6 marks. The key point is: do not get stuck on one question.

  \item A good strategy is to first answer the questions you are most confident about before
  attempting the others. In other words, answering questions strictly serially during an
  exam can be suboptimal.

  \item Do not waste time answering an unrelated question in the hope of gaining partial marks.
  Marking will be based on the question that was actually asked.

  \item Where necessary, justify your answers with clear and concise explanations, without being
  overly wordy.
\end{itemize}

\newpage

\questiontitle{1}{6}

Let $S[1 \ldots m]$ and $T[1 \ldots n]$ be two strings. An \emph{overlap} of
$S$ into $T$ is an integer $\ell$ such that the suffix
$S[m-\ell+1 \ldots m]$ is exactly equal to the prefix $T[1 \ldots \ell]$.

Write pseudocode for an efficient algorithm that accepts $S$ and $T$ and
outputs the maximum overlap length $\ell$.

For this problem, you may assume that \texttt{zalg} is available as a built-in
function that takes a string as input and returns the array of its Z-values. Do
not write pseudocode for \texttt{zalg}. State the worst-case time and space
complexity of your algorithm.

\marksbox{6}

\workingpage
\newpage

\questiontitle{2}{6}

In the Boyer--Moore algorithm, the extended bad-character rule stores
$R_k(x)$, the position of the rightmost occurrence of character $x$ strictly to
the left of position $k$ in the pattern $pat[1 \ldots m]$, or $0$ if no such
occurrence exists.

\begin{enumerate}
  \item[(i)] Write pseudocode to construct the full $m \times |\Sigma|$ table
  of $R_k(x)$ values for a fixed alphabet $\Sigma$.

  \item[(ii)] Suppose a right-to-left scan mismatches at pattern position $k$
  against text character $x$. Explain how the extended bad-character shift is
  computed from $R_k(x)$, and state the preprocessing time and space
  complexity.
\end{enumerate}

\marksbox{6}

\workingpage
\newpage

\questiontitle{3}{6=2+4}

Assume that $\$$ is a unique terminal symbol that is lexicographically smaller
than every other character. The string $S=\texttt{banana\$}$ has
\[
  L = BWT(S)=\texttt{annb\$aa}
\]
and suffix array
\[
  SA = [7,6,4,2,1,5,3].
\]

Using $L$ as a Burrows--Wheeler search index, search for the pattern
$pat=\texttt{ana}$ by backward search.

\begin{enumerate}
  \item[(i)] Construct the needed $rank(c)$ values and the occurrence counts
  needed during the search.

  \item[(ii)] Show the values of $sp$ and $ep$ after each processed pattern
  character. Then report the multiplicity of $pat$ in $S$ and the corresponding
  starting positions using the suffix array.
\end{enumerate}

\marksbox{6}

\workingpage
\newpage

\questiontitle{4}{6=3+3}

These questions relate to Ukkonen's linear-time suffix tree construction
algorithm. Let $I_n$ be the implicit suffix tree for a string $S[1 \ldots n]$
that does not contain the terminal character $\$$.

\begin{enumerate}
  \item[(i)] In the final phase that appends the unique terminal character
  $\$$, can any explicit suffix extension be completed by Rule 3? Answer Yes or
  No and justify your answer.

  \item[(ii)] Explain why appending $\$$ converts the final implicit suffix tree
  into a regular suffix tree with one leaf for every suffix. In your answer,
  connect this to the role of Rule 2 and the rapid leaf-extension trick.
\end{enumerate}

\marksbox{6}

\workingpage
\newpage

\questiontitle{5}{6=3+3}

\begin{enumerate}
  \item[(i)] Decode the following Elias--Omega codeword to recover the positive
  integer $N$ that it encodes:
  \[
    \texttt{00000010100}
  \]
  Show each component read during decoding.

  \item[(ii)] The following tuples were produced by the LZ77 method:
  \[
    \langle 0,0,m\rangle,\quad
    \langle 0,0,i\rangle,\quad
    \langle 0,0,s\rangle,\quad
    \langle 3,3,p\rangle,\quad
    \langle 4,4,!\rangle.
  \]
  Decode these tuples to regenerate the original string.
\end{enumerate}

\marksbox{6}

\workingpage
\newpage

\questiontitle{6}{6}

Perform one iteration of the Miller--Rabin primality test on $n = 91$ using
base $a = 3$.

Your answer should:
\begin{itemize}
  \item express $n-1$ in the form $2^s t$, where $t$ is odd;
  \item compute the sequence of values checked by this Miller--Rabin iteration;
  \item state whether this base proves that $91$ is composite, or whether this
  iteration is inconclusive.
\end{itemize}

\marksbox{6}

\workingpage
\newpage

\questiontitle{7}{6}

In a Fibonacci heap, after \textsc{Extract-Min} removes $H.min$, the children
of the removed root are promoted into the root list. The consolidation step
then repeatedly links roots of equal degree until no two roots in the root list
have the same degree.

Write pseudocode for this consolidation step. Your pseudocode should use an
auxiliary array indexed by degree and should link two roots of equal degree by
making the larger-key root a child of the smaller-key root.

State the worst-case actual time and amortized time of \textsc{Extract-Min}.

\marksbox{6}

\workingpage
\newpage

\questiontitle{8}{6}

A dynamic array starts empty. The first append creates capacity $1$. Whenever
an append is attempted on a full array, a new array with twice the previous
capacity is allocated, all existing elements are copied into it, and then the
new element is appended.

Using the potential method, determine the amortized time complexity of an
append operation over a sequence of $n$ appends. Use the potential function
\[
  \Phi(A_i)=2\cdot size(A_i)-capacity(A_i).
\]

Verify the required boundary conditions for $\Phi$, and show the amortized
cost calculation for both a non-resizing append and a resizing append.

\marksbox{6}

\workingpage
\newpage

\questiontitle{9}{6}

Let $P$ be a node in a B-tree of minimum degree $t$. The node stores sorted keys
$key_1,\ldots,key_r$ and, if it is not a leaf, child pointers
$child_0,\ldots,child_r$.

Write pseudocode for a function \textsc{Search}$(P,x)$ that returns whether key
$x$ occurs in the subtree rooted at $P$. Your pseudocode should use binary
search inside each visited node before deciding which child subtree to search.

Justify that your procedure terminates correctly, and state its worst-case time
complexity in terms of the B-tree height $h$ and minimum degree $t$.

\marksbox{6}

\workingpage
\newpage

\questiontitle{10}{6}

Use the Hungarian algorithm to solve the following minimum-cost assignment
problem.

\[
\begin{array}{c|ccc}
      & \text{Task 1} & \text{Task 2} & \text{Task 3} \\
\hline
\text{Agent 1} & 9 & 2 & 7 \\
\text{Agent 2} & 6 & 4 & 3 \\
\text{Agent 3} & 5 & 8 & 1
\end{array}
\]

Show the row and column reductions you perform, identify a minimum-cost
perfect matching, and state the total assignment cost.

\marksbox{6}

\workingpage

\vfill
\begin{center}
  End of Exam
\end{center}

\end{document}
